给服务集群写了条自动扩容规则：CPU 利用率超过 \(70\%\) 就加机器，低于 \(30\%\) 就撤。上线之后机器数每隔几分钟就在两个值之间抖一次。

规则本身没写错——CPU 高了确实该加机器。错的是这条规则和它所管的系统合起来构成了一个回路，而回路的行为读不出来。新机器要几分钟才真正接管流量，在这段时间里控制器一直在对已经过时的读数做反应，于是加多了；等负载降下去又撤多了。增益、时滞、震荡——这三个词是同一件事的三个侧面，而它们没有一个能从``CPU 超过 70\% 就加机器''这句话里看出来。

一旦你的动作会改变你下一次观测到的东西，``这条规则对不对''就不再是一个可以单独回答的问题。要回答它，必须把系统、传感器、控制器、执行这四样东西当成一个整体来分析——这就是本部分的入口。

前一部分只是旁观系统自己往哪走，这里要往里面加一只手。加进去之后，问题从``系统会去哪''变成``怎样让它去我们要它去的地方''，而``我们要它去哪''本身必须先被写清楚：是别偏离目标太远，是用最少的燃料到达，还是在有限时间内不越过某条红线——不同的写法通向完全不同的数学。

\begin{center}
\begin{tikzpicture}[x=1cm,y=1cm,font=\small,>=stealth]
  % --- 闭环回路 ---
  \draw[black!50,fill=blue!8] (0,2.4) rectangle (2.3,3.3);
  \node at (1.15,2.85) {系统};
  \draw[black!50,fill=blue!8] (3.3,2.4) rectangle (5.5,3.3);
  \node at (4.4,2.85) {传感器};
  \draw[black!50,fill=green!9] (3.3,0.9) rectangle (5.5,1.8);
  \node at (4.4,1.35) {估计器};
  \draw[black!50,fill=green!9] (0,0.9) rectangle (2.3,1.8);
  \node at (1.15,1.35) {控制器};

  \draw[->,black!60] (2.35,2.85) -- (3.25,2.85);
  \node[font=\footnotesize,black!70] at (2.8,3.12) {\(x\)};
  \draw[->,black!60] (4.4,2.35) -- (4.4,1.85);
  \node[font=\footnotesize,black!70] at (4.72,2.1) {\(y\)};
  \draw[->,black!60] (3.25,1.35) -- (2.35,1.35);
  \node[font=\footnotesize,black!70] at (2.8,1.62) {\(\hat x\)};
  \draw[->,black!60] (1.15,1.85) -- (1.15,2.35);
  \node[font=\footnotesize,black!70] at (1.47,2.1) {\(u\)};

  \draw[->,black!40] (1.15,3.75) -- (1.15,3.35);
  \node[font=\footnotesize,black!60] at (1.15,3.95) {过程噪声};
  \draw[->,black!40] (4.4,3.75) -- (4.4,3.35);
  \node[font=\footnotesize,black!60] at (4.4,3.95) {观测噪声};

  % --- 三级放松 ---
  \draw[black!45,fill=black!4] (6.3,3.0) rectangle (12.7,3.9);
  \node[align=left,anchor=west,font=\footnotesize] at (6.5,3.45)
    {模型已知、状态可测：直接解最优性条件};
  \draw[black!45,fill=black!4] (6.3,1.85) rectangle (12.7,2.75);
  \node[align=left,anchor=west,font=\footnotesize] at (6.5,2.3)
    {模型已知、状态看不全：先估计，再把估计当状态用};
  \draw[black!45,fill=black!4] (6.3,0.7) rectangle (12.7,1.6);
  \node[align=left,anchor=west,font=\footnotesize] at (6.5,1.15)
    {模型未知：只能采样，逼近同一个不动点};
  \draw[->,black!45] (5.6,3.0) -- (6.25,3.45);
  \draw[->,black!45] (5.6,2.3) -- (6.25,2.3);
  \draw[->,black!45] (5.6,1.6) -- (6.25,1.15);
\end{tikzpicture}

\emph{这一部分做的事，是把这个回路上的假设一条条拿掉。每拿掉一条，能给出的保证就弱一层，需要的数据就多一批——而 Bellman 那个递归始终没有换。}
\end{center}

\subsection{五条贯穿始终的判断}

\textbf{开环与闭环的差别不在精度，在有没有把误差喂回来。}开环控制器照预定计划执行，任何没被建模的扰动都会一路累积下去；闭环用当前的偏差决定下一步，扰动因此被持续修正。但这样就引入了一个回路，而回路会震荡：增益越大纠正越快，也越容易过冲，再叠上执行的时滞就成了开头那种抖动。反馈买来的是抗扰动能力，付出的是稳定性从此需要单独论证。

\textbf{能不能控制、能不能看清，是两个独立的问题，而且都是结构性的。}可控性问的是输入能否把状态推到任意目标，可观测性问的是从输出能否反推出状态。这两件事由系统的结构决定，调参解决不了——某个方向本来就不可控，再大的输入也推不动；不可观测的分量，再好的滤波器也估不出来。它们该在设计控制器之前诊断，而不是在效果不好之后才回头查。

\textbf{最优性有两种写法，短处刚好相反。}最大值原理给出的是最优轨线必须满足的条件，结果是一个两点边值问题——算一条轨线便宜，但换个初值要重算。HJB 给出的是最优策略本身，一个覆盖所有状态的函数——一次求解处处可用，但要在整个状态空间上求解，于是撞上维数灾难。这个取舍在后半部分原样重现：策略梯度直接优化一条策略，值迭代试图把所有状态一起管了。

\textbf{先估计再控制是一个定理，不是常识。}线性动力学、Gaussian 噪声、二次代价，这三条同时成立时分离定理才成立：先算最优状态估计，再把估计当成真状态代进确定性最优控制律，合起来仍然是最优的。缺任何一条，这个拆分就退化成一个需要单独验证的工程近似——它可能仍然好用，但``最优''两个字已经不能再说了。

\textbf{模型未知时，所有保证都要重新对账。}Bellman 不动点还在，但方程里的转移概率从已知变成只能采样，于是收敛性依赖的东西整个换了一套：压缩模数让位给步长条件，而当自举、函数逼近、离策略这三样同时出现时，连收敛本身都可能没有。这不是哪个算法不够好——三个各自安全的选择凑在一起就不安全了，而它们各自的收敛性证明恰好都用到了另外两个不成立的东西。

\subsection{两章怎么安排}

\hyperref[ch:135]{``最优控制：从反馈稳定到全局最优''}处理模型已知的情形。头三篇是主干：从开环到闭环建立反馈的必要性，最大值原理与 HJB 给出前面说的那两种最优性写法。中间两篇处理``看不全状态''——Kalman 滤波用预测与更新的循环从噪声里估计状态，LQG 则说明什么条件下可以把它和最优反馈分开设计。最后两篇偏工程：DDP 用局部二次模型改进整条轨迹，系统辨识从输入输出数据估计动力学，其中``数据有没有激励关键方向''这个问题要排在拟合之前——激励不足时参数估计会安静地失败。

\hyperref[ch:136]{``强化学习：模型未知时怎样逼近最优决策''}处理连模型都没有的情形。前四篇把地基铺完：MDP 把决策问题写成随机过程，Bellman 方程与值迭代给出不动点视角，Monte Carlo 与随机逼近说明为什么用样本代替期望仍然收敛，Q 学习则是``把期望换成样本、最大值留在原地''这一步的直接产物。中间三篇是全章最要紧的部分——线性价值函数逼近解的其实是投影 Bellman 方程，它的不动点不是最佳逼近；Deadly Triad 说明三个安全选择凑在一起为什么会发散；DQN 的慢目标与回放则是对这个耦合的工程缓解，明说是缓解而不是解决。最后三篇转到策略一侧：离策略校正用权重交换偏差与方差，策略梯度绕开了取最大值这个子问题，Actor--Critic 在两者之间搭桥。

两章共享 Bellman 的递归，但不要压缩成``已知模型解方程、未知模型逼近同一个不动点''。策略梯度就不走这条路，它直接优化动作分布；部分可观测、平均回报、纯离线数据还会各自改变问题的形式。更稳的主线是先把四件事说清楚——控制器看得见什么，能改变什么，在优化什么，数据由谁产生——之后再挑方法。

前面需要的底子都已经备好。反馈为什么能稳住系统，用的是\hyperref[sec:15-Differential-Equations-and-Dynamical-Systems/02-Stability-and-Dynamical-Systems/02]{``稳定性与 Lyapunov 函数''}那套能量思想；MDP 去掉动作就退回\hyperref[sec:06-Random-Processes/03-Markov-Chains/01]{``Markov 性质、转移矩阵与多步演化''}的链；最大值原理导出的两点边值问题，计算上交给\hyperref[sec:08-Numerical-Analysis/06-ODE/04]{``边值问题：从打靶到全局离散''}。反过来，这一部分给出的闭环视角在别处也用得上：任何一个把自己的输出喂回输入的系统——推荐、风控、定价——都已经不是纯粹的预测问题了，\hyperref[ch:130]{``微分方程与动力系统：从局部变化规律到整体行为''}里关于长期行为的那些判断在那里同样适用。
